\chapter{Tableau de pilotage des 55 séances}
\small
\begin{longtable}{cclp{.42\textwidth}}
\toprule
N° & Date & Chapitre & Objet\\
\midrule\endhead
1 & 15/09/2026 & Suites numériques & Suite arithmétique : reconnaître et modéliser\\
2 & 17/09/2026 & Suites numériques & Suite arithmétique : terme général\\
3 & 22/09/2026 & Suites numériques & Somme des termes d'une suite arithmétique\\
4 & 24/09/2026 & Suites numériques & Suite géométrique : reconnaître et modéliser\\
5 & 29/09/2026 & Suites numériques & Suite géométrique : terme général et somme\\
6 & 01/10/2026 & Suites numériques & Bilan et évaluation formative\\
7 & 06/10/2026 & Fonctions exponentielles & Définition et sens de variation de $a^x$\\
8 & 08/10/2026 & Fonctions exponentielles & Puissances et propriétés algébriques\\
9 & 13/10/2026 & Fonctions exponentielles & Fonctions $x\mapsto k a^x$ : variations et graphique\\
10 & 15/10/2026 & Fonctions exponentielles & Taux d'évolution moyen\\
11 & 03/11/2026 & Fonctions exponentielles & Lecture de représentations graphiques\\
12 & 05/11/2026 & Fonctions exponentielles & Équations exponentielles simples\\
13 & 10/11/2026 & Fonctions exponentielles & Modélisation par une fonction exponentielle\\
14 & 12/11/2026 & Fonctions exponentielles & Bilan et évaluation du chapitre\\
15 & 17/11/2026 & Fonction logarithme décimal & Définition du logarithme décimal\\
16 & 19/11/2026 & Fonction logarithme décimal & Sens de variation et représentation\\
17 & 24/11/2026 & Fonction logarithme décimal & Propriétés algébriques\\
18 & 26/11/2026 & Fonction logarithme décimal & Résoudre $10^x=b$\\
19 & 01/12/2026 & Fonction logarithme décimal & Résoudre $a^x=b$\\
20 & 03/12/2026 & Fonction logarithme décimal & Équations $x^a=b$ et inégalités simples\\
21 & 08/12/2026 & Fonction logarithme décimal & Problèmes contextualisés\\
22 & 10/12/2026 & Fonction logarithme décimal & Bilan et évaluation\\
23 & 15/12/2026 & Fonction inverse & Définition et courbe de la fonction inverse\\
24 & 17/12/2026 & Fonction inverse & Comportement aux bornes et asymptotes\\
25 & 05/01/2027 & Fonction inverse & Dérivée et sens de variation\\
26 & 07/01/2027 & Fonction inverse & Combinaisons avec des fonctions polynomiales\\
27 & 12/01/2027 & Fonction inverse & Signe et factorisation\\
28 & 14/01/2027 & Fonction inverse & Dérivées de polynômes de degré au plus 3\\
29 & 19/01/2027 & Fonction inverse & Étude complète et représentation\\
30 & 21/01/2027 & Fonction inverse & Bilan et évaluation\\
31 & 26/01/2027 & Statistiques à deux variables & Nuage de points\\
32 & 28/01/2027 & Statistiques à deux variables & Lecture et interprétation du nuage\\
33 & 02/02/2027 & Statistiques à deux variables & Ajustement affine\\
34 & 04/02/2027 & Statistiques à deux variables & Interpolation\\
35 & 09/02/2027 & Statistiques à deux variables & Extrapolation\\
36 & 11/02/2027 & Statistiques à deux variables & Changement de variable\\
37 & 02/03/2027 & Statistiques à deux variables & Choix du modèle et ajustement\\
38 & 04/03/2027 & Statistiques à deux variables & Bilan et évaluation\\
39 & 09/03/2027 & Probabilités conditionnelles & Conditionnement\\
40 & 11/03/2027 & Probabilités conditionnelles & Arbres pondérés\\
41 & 16/03/2027 & Probabilités conditionnelles & Produit des probabilités le long d'un chemin\\
42 & 18/03/2027 & Probabilités conditionnelles & Indépendance\\
43 & 23/03/2027 & Probabilités conditionnelles & Partition et probabilités totales\\
44 & 25/03/2027 & Probabilités conditionnelles & Problèmes contextualisés\\
45 & 30/03/2027 & Probabilités conditionnelles & Entraînement mixte\\
46 & 01/04/2027 & Probabilités conditionnelles & Bilan et évaluation\\
47 & 06/04/2027 & Variables aléatoires discrètes & Espérance d'une variable aléatoire\\
48 & 08/04/2027 & Variables aléatoires discrètes & Loi binomiale : reconnaître une situation\\
49 & 27/04/2027 & Variables aléatoires discrètes & Triangle de Pascal et coefficients binomiaux\\
50 & 29/04/2027 & Variables aléatoires discrètes & Probabilités binomiales\\
51 & 04/05/2027 & Variables aléatoires discrètes & Bilan du chapitre et évaluation\\
52 & 11/05/2027 & Révisions et consolidation & Diagnostic des acquis\\
53 & 13/05/2027 & Révisions et consolidation & Réactivation : fonctions exponentielles\\
54 & 18/05/2027 & Révisions et consolidation & Réactivation : logarithme décimal\\
55 & 20/05/2027 & Révisions et consolidation & Réactivation : fonction inverse\\
\bottomrule
\end{longtable}
\normalsize
